% ============================================================================
% 英文摘要
% 写作要点：
%   1. 与中文摘要严格三段对应
%   2. 全文使用 Times New Roman 12pt（小四）
%   3. Keywords 使用英文分号分隔，与中文关键词逐项对应
% ============================================================================

\begin{abstracten}

\noindent{\fontspec{Times New Roman}\bfseries\zihao{-4} Abstract: }Medical inquiry training depends on clinical cases, standardized patients and timely teacher feedback. In routine university teaching, however, these resources are not always available for repeated practice. Large language models make virtual patients more feasible, but a prompt-only solution may easily cross role boundaries, fabricate examination results, or provide vague feedback. This thesis therefore designs and implements a virtual consultation system for medical inquiry training, with an emphasis on interactive practice, controllable case facts and reviewable training results.

The system follows a front-end / back-end separated architecture. The front end is developed with Vue 3, TypeScript and Element Plus, while the back end uses Spring Boot 3.3.6, Java 17 and MyBatis-Plus. MySQL stores users, cases and training records; Redis keeps verification codes, login locks and dialogue context; MinIO manages medical images and attachments; Docker Compose is used for deployment. For students, teachers and administrators, the system provides AI consultation training, case-library management, learning analytics and system administration, together with JWT authentication, RBAC authorization, ECharts visualization and SSE streaming.

For the AI workflow, the system separates a single prompt into three roles: PatientAgent, DiagnosisAgent and FeedbackAgent. AgentRouter selects the proper role according to the training stage. A keyword-based RAG module supplies medical knowledge, Redis memory keeps recent dialogue turns, and clinical tools return physical examination, auxiliary examination and imaging results directly from case data instead of allowing the model to invent objective findings. After training, the system outputs structured JSON scores for inquiry, diagnosis, treatment and communication, with rule-based fallback when parsing fails. Deployment and functional, performance and security tests show that the system can support a complete medical inquiry training loop.

\keywordsen{Medical Education; Virtual Consultation; Large Language Model; AI Agent; RAG; Front-end / Back-end Separation}

\end{abstracten}
